\documentclass[11pt,letterpaper]{article}

% --- Packages ---
\usepackage[margin=1in]{geometry}
\usepackage{graphicx}
\usepackage{booktabs}
\usepackage{longtable}
\usepackage{tabularx}
\usepackage{enumitem}
\usepackage{hyperref}
\usepackage{xcolor}
\usepackage{fancyhdr}
\usepackage{listings}
\usepackage{titlesec}
\usepackage[T1]{fontenc}
\usepackage{lmodern}
\usepackage{parskip}
\usepackage{array}
\usepackage{microtype}

% --- Colours ---
\definecolor{codebg}{gray}{0.95}
\definecolor{codeframe}{gray}{0.75}
\definecolor{linkblue}{HTML}{0066CC}
\definecolor{cautionred}{HTML}{CC0000}
\definecolor{notegray}{HTML}{555555}

% --- Hyperlink styling ---
\hypersetup{
    colorlinks=true,
    linkcolor=linkblue,
    urlcolor=linkblue,
    citecolor=linkblue,
}

% --- Code listing style ---
\lstset{
    backgroundcolor=\color{codebg},
    frame=single,
    rulecolor=\color{codeframe},
    basicstyle=\ttfamily\small,
    breaklines=true,
    breakatwhitespace=false,
    columns=fullflexible,
    keepspaces=true,
    showstringspaces=false,
    xleftmargin=4pt,
    xrightmargin=4pt,
    aboveskip=8pt,
    belowskip=8pt,
}

% --- Header / Footer ---
\pagestyle{fancy}
\fancyhf{}
\fancyhead[L]{\textbf{X-CAVATE Original Refactoring Report}}
\fancyhead[R]{\thepage}
\renewcommand{\headrulewidth}{0.4pt}

% --- Custom commands ---
\newcommand{\code}[1]{\texttt{#1}}
\tolerance=400
\emergencystretch=2em
\newcolumntype{L}[1]{>{\raggedright\arraybackslash}p{#1}}
\newcolumntype{C}[1]{>{\centering\arraybackslash}p{#1}}
\newcolumntype{R}[1]{>{\raggedleft\arraybackslash}p{#1}}

% --- Title ---
\title{
    \vspace{-1cm}
    \textbf{X-CAVATE Original Refactoring Report}\\[6pt]
    \large Detailed Line-by-Line Mapping:\\
    \code{xcavate\_11\_30\_25.py} (5,568 lines) $\to$ \code{xcavate/} Package (6,201 lines)
}
\author{}
\date{March 8, 2026}

\begin{document}

\maketitle
\thispagestyle{fancy}

\tableofcontents
\newpage

%=============================================================================
\section{Overview}
%=============================================================================

The original X-CAVATE was a single monolithic Python script with 3 user-defined functions and 30+ global variables. The refactored version is a modular package with 25 classes, $\sim$95 functions, dataclass configuration, abstract base classes, and spatial indexing. This document traces every section of the original code to its new location, documents what changed algorithmically, and quantifies the reductions.

\subsection{Size Comparison}

\begin{table}[h]
\centering
\begin{tabular}{l r r r}
\toprule
\textbf{Metric} & \textbf{Original} & \textbf{Refactored} & \textbf{Change} \\
\midrule
Total lines               & 5,568  & 6,201  & +633 (+11.4\%) \\
Files                     & 1      & 27     & +26 \\
Functions/methods         & 3      & $\sim$95 & +92 \\
Classes                   & 0      & 25     & +25 \\
Global variables          & 30+    & 0      & $-$30 \\
G-code blocks (copy-pasted) & 6    & 3 subclasses + 1 ABC & $-$3 dup.\ blocks \\
Plot blocks (copy-pasted) & 5--7   & 2 parameterized funcs & $-$3--5 blocks \\
Coordinate output blocks  & 2 (SM + MM) & 1 parameterized func & $-$1 block \\
\bottomrule
\end{tabular}
\caption{Overall size comparison.}
\end{table}

%=============================================================================
\section{Section-by-Section Mapping}
%=============================================================================

%-----------------------------------------------------------------------------
\subsection{Imports and Setup}
%-----------------------------------------------------------------------------

\textbf{Original:} Lines 1--37 (37 lines)\\
\textbf{Refactored:} \code{xcavate/\_\_init\_\_.py} (3 lines), \code{\_\_main\_\_.py} (6 lines) --- total 9 lines

Imports pandas, numpy, argparse, matplotlib, plotly; creates output directories.

\textbf{What changed:} Output directory creation moved to \code{XcavateConfig.ensure\_output\_dirs()} (4 lines in \code{config.py}). Imports are now per-module (each file imports only what it needs). Random seed removed (not used anywhere after initial set).

%-----------------------------------------------------------------------------
\subsection{Argument Parsing and Global Variables}
%-----------------------------------------------------------------------------

\textbf{Original:} Lines 38--174 (137 lines)\\
\textbf{Refactored:} \code{cli.py} (192 lines) + \code{config.py} (131 lines) --- total 323 lines

\textbf{What changed:}

\begin{enumerate}
    \item \textbf{30+ global variables replaced by \code{XcavateConfig} dataclass.} Every parameter is now a typed field with a default value and docstring, rather than a bare \code{variable = args.variable} assignment.

    \item \textbf{Custom G-code file paths.} The original hardcoded paths like \code{startExtrusionCode = 'inputs/custom/start\_extrusion\_code.txt'} (lines 136--161). The refactored code uses \code{CustomCodes.load\_from\_dir(path)} which reads all 12 template files once from a configurable directory.

    \item \textbf{New parameters added:}
    \begin{itemize}
        \item \code{--algorithm} (dfs or sweep\_line) for pluggable pathfinding
        \item \code{--overlap\_algorithm} (retrace or consecutive) for selectable overlap
        \item \code{--output\_dir} (configurable output location, was hardcoded to \code{outputs/})
    \end{itemize}

    \item \textbf{Type safety.} Enums for \code{PrinterType}, \code{PathfindingAlgorithm}, \code{OverlapAlgorithm} replace raw integers and strings.
\end{enumerate}

Example: original global variable declarations (lines 101--174):

\begin{lstlisting}[language=Python]
multimaterial = args.multimaterial    # -> config.multimaterial
tolerance = args.tolerance            # -> config.tolerance
nozzle_OD = args.nozzle_diameter      # -> config.nozzle_diameter
nozzle_radius = nozzle_OD/2           # -> config.nozzle_radius (property)
scaleFactor = args.scale_factor       # -> config.scale_factor
printhead1_axis = 'A'                 # -> config.axis_1 (configurable)
startExtrusionCode = 'inputs/...'     # -> CustomCodes.load_from_dir()
\end{lstlisting}

%-----------------------------------------------------------------------------
\subsection{File Reading and Preprocessing}
%-----------------------------------------------------------------------------

\textbf{Original:} Lines 177--324 (148 lines)\\
\textbf{Refactored:} \code{io/reader.py} (298 lines)

Reads network file, removes blank lines/headers, writes preprocessed.txt/inlets.txt/outlets.txt, reads into pandas DataFrame, converts to numpy, scales, rounds, matches inlet/outlet nodes.

\subsubsection{Detailed Breakdown}

\begin{table}[h]
\centering
\footnotesize
\begin{tabularx}{\textwidth}{L{1.5cm} X L{4cm} R{0.8cm}}
\toprule
\textbf{Old Lines} & \textbf{What It Does} & \textbf{Refactored Location} & \textbf{Lines} \\
\midrule
177--205 & Strip blank lines, extract vessel headers, write \code{preprocessed.txt} & \code{read\_network\_file()} \newline reader.py:29--89 & 61 \\
206--234 & Parse inlet/outlet file, write \code{inlets.txt}, \code{outlets.txt} & \code{read\_inlet\_outlet\_\allowbreak file()} \newline reader.py:92--157 & 66 \\
238--268 & Read into pandas, detect columns, convert to numpy & Merged into \code{read\_network\_file()} & incl. \\
279--305 & Unit conversion (cm$\to$mm), scaling, rounding & \code{preprocess\_coordinates()} \newline reader.py:172--226 & 55 \\
307--324 & Match inlet/outlet coords to nodes (exact \code{==}) & \code{match\_inlet\_outlet\_\allowbreak nodes()} \newline reader.py:229--298 & 70 \\
\bottomrule
\end{tabularx}
\caption{File reading detailed breakdown.}
\end{table}

\subsubsection{Key Algorithmic Changes}

\begin{enumerate}
    \item \textbf{No intermediate files.} The original wrote \code{outputs/graph/preprocessed.txt}, \code{inlets.txt}, and \code{outlets.txt} as intermediate steps. The refactored code processes everything in-memory, eliminating 3 file writes and 3 file reads.

    \item \textbf{Inlet/outlet matching: exact float \code{==} vs KD-tree nearest-neighbor.}
    \begin{itemize}
        \item Original (lines 307--324): Triple-nested loop comparing \code{x == x}, \code{y == y}, \code{z == z} individually. Brittle with floating-point values.
        \item Refactored: \code{scipy.spatial.cKDTree} with \code{.query()} for $O(\log N)$ nearest-neighbor. Warns if distance $> 5.0$.
    \end{itemize}

    \item \textbf{Coordinate rounding: triple-nested Python loop vs vectorized numpy.}
    \begin{itemize}
        \item Original (lines 292--305): \code{for i in range(len(points)):} $\to$ $O(N \times C)$ Python-level iteration.
        \item Refactored: \code{np.round(points, decimals=num\_decimals)} --- single vectorized call.
    \end{itemize}

    \item \textbf{Two-pass parsing.} The refactored reader scans the file twice: first to count vessels and points (for pre-allocation), then to fill the array. The original used Python list appending and pandas conversion.
\end{enumerate}

%-----------------------------------------------------------------------------
\subsection{Original Network Plot}
%-----------------------------------------------------------------------------

\textbf{Original:} Lines 328--394 (67 lines)\\
\textbf{Refactored:} \code{viz/plotting.py:77-136} (\code{create\_original\_network\_plot}, 60 lines)

Creates 3D Plotly scatter plot of pre-interpolation network, one trace per vessel, with slider. The function was extracted and parameterized. The original inline block was nearly identical to 4 other plotting blocks; the refactored version is a single reusable function.

\textbf{Note:} In the initial refactoring, \code{create\_original\_network\_plot()} existed but was \textbf{never called} (dead code). This was fixed in the parity update by adding the call in \code{\_write\_plots()}.

%-----------------------------------------------------------------------------
\subsection{Interpolation}
%-----------------------------------------------------------------------------

\textbf{Original:} Lines 396--487 (92 lines)\\
\textbf{Refactored:} \code{core/preprocessing.py:45-181} (\code{interpolate\_network}, 137 lines)

Densifies network so consecutive point spacing does not exceed nozzle radius.

\subsubsection{Key Algorithmic Changes}

\begin{enumerate}
    \item \textbf{$O(N)$ single \code{np.concatenate} vs repeated \code{np.insert}.}
    \begin{itemize}
        \item Original (lines 437--481): For each gap, calls \code{np.insert()} which copies the entire array. For a network gaining 30,000 points, this is $O(N^2)$ total copy operations.
        \item Refactored: Builds a list of segment arrays, then calls \code{np.concatenate()} once at the end. $O(N)$ total.
    \end{itemize}

    \item \textbf{Flag column encoding.}
    \begin{itemize}
        \item Original: Sets column 4 (or 5) to \code{500} for vessel start points, reuses the same column.
        \item Refactored: Appends a dedicated flag column: \code{500} = vessel start, \code{400} = interpolated, \code{0} = original. Cleaner separation.
    \end{itemize}

    \item \textbf{Line count increase (+45 lines):} Due to more comprehensive docstrings, column-count handling, and the \code{get\_vessel\_boundaries()} helper function (68 lines).
\end{enumerate}

%-----------------------------------------------------------------------------
\subsection{Graph Construction (Endpoints, Branchpoints, Adjacency)}
%-----------------------------------------------------------------------------

\textbf{Original:} Lines 489--906 (418 lines)\\
\textbf{Refactored:} \code{core/graph.py} (836 lines)

Finds endpoints, detects branchpoints via distance search, validates daughter pairs via parametric line test, builds adjacency graph.

\subsubsection{Detailed Sub-section Mapping}

\begin{table}[h]
\centering
\footnotesize
\begin{tabularx}{\textwidth}{L{1.6cm} X L{4cm} R{0.8cm}}
\toprule
\textbf{Old Lines} & \textbf{Step} & \textbf{Refactored Function} & \textbf{Lines} \\
\midrule
489--540 (52) & Find endpoints per vessel & \code{\_find\_endpoints()} \newline graph.py:162--203 & 42 \\
541--559 (19) & Re-match inlets/outlets post-interpolation & Handled in pipeline & 0 \\
560--742 (183) & Find branchpoints (distance search + line test) & \code{\_find\_branchpoints()} + \newline \code{\_validate\_daughter\_pairs()} \newline graph.py:210--540 & 331 \\
828--906 (79) & Build adjacency graph + remove cross-edges & \code{\_build\_adjacency()} + \newline \code{\_remove\_daughter\_\allowbreak cross\_edges()} \newline graph.py:596--750 & 155 \\
(inline) & Write \code{graph.txt}, \code{special\_nodes.txt} & \code{write\_graph()} + \code{write\_special\_nodes()} \newline graph.py:757--835 & 79 \\
\bottomrule
\end{tabularx}
\caption{Graph construction sub-section mapping.}
\end{table}

\subsubsection{Key Algorithmic Changes}

\begin{enumerate}
    \item \textbf{Branchpoint detection: $O(N \times E)$ nested loops vs $O(N)$ KD-tree.}
    \begin{itemize}
        \item Original (lines 575--625): For each endpoint, iterates through ALL points to find two nearest points outside the endpoint's vessel. $O(E \times N)$.
        \item Refactored: Builds \code{cKDTree} once, queries \code{.query(point, k=10)} for $k$ nearest neighbors, filters by vessel. $O(E \times k \times \log N)$.
    \end{itemize}

    \item \textbf{Parametric line test preserved.} The daughter validation logic is preserved but extracted into \code{\_select\_pair\_by\_line\_test()}.

    \item \textbf{Line count doubled (+418 lines):} Due to 12 separate functions (vs inline code), comprehensive docstrings, and additional helper functions for edge cases (e.g., \code{\_brute\_force\_nearest()} fallback).
\end{enumerate}

%-----------------------------------------------------------------------------
\subsection{DFS Pathfinding}
%-----------------------------------------------------------------------------

\textbf{Original:} Lines 911--995 (85 lines)\\
\textbf{Refactored:} \code{core/pathfinding.py} (620 lines)

Recursive DFS with collision checking to generate print passes --- refactored into iterative DFS with KD-tree collision checking; also includes sweep-line alternative.

\subsubsection{Detailed Sub-section Mapping}

\begin{table}[h]
\centering
\footnotesize
\begin{tabularx}{\textwidth}{L{1.6cm} X L{4cm} R{0.8cm}}
\toprule
\textbf{Old Lines} & \textbf{Function} & \textbf{Refactored Location} & \textbf{Lines} \\
\midrule
911--935 (25) & \code{is\_valid()} --- checks if printing a node blocks others below & \code{CollisionDetector.\allowbreak is\_valid()} \newline pathfinding.py:110--155 & 46 \\
937--950 (14) & \code{find\_lowest\_unvisited()} --- $O(N)$ linear scan & \code{\_pop\_lowest\_unvisited()} \newline with min-heap \newline pathfinding.py:326--348 & 23 \\
953--982 (30) & \code{dfs()} --- recursive DFS & \code{iterative\_dfs()} with stack \newline pathfinding.py:351--407 & 57 \\
983--995 (13) & Main loop calling DFS per pass & \code{DFSStrategy.\allowbreak generate\_print\_passes()} \newline pathfinding.py:265--323 & 59 \\
(new) & Sweep-line alternative & \code{SweepLineStrategy} \newline pathfinding.py:414--560 & 147 \\
(new) & Collision detector class & \code{CollisionDetector} \newline pathfinding.py:48--200 & 153 \\
\bottomrule
\end{tabularx}
\caption{DFS pathfinding sub-section mapping.}
\end{table}

\subsubsection{Key Algorithmic Changes}

\begin{enumerate}
    \item \textbf{Recursive $\to$ Iterative DFS.} Original: Python recursive function that hits \code{RecursionError} on networks with $>$1,000 nodes after interpolation. Our test network (31,271 points) crashes at default limit. Refactored: Explicit stack-based DFS, no recursion limit.

    \item \textbf{\code{is\_valid()} collision check: $O(N)$ vs $O(k \log N)$.} Original: computes XY distance to ALL unvisited nodes. Refactored: \code{cKDTree.query\_ball\_point(xy, nozzle\_radius)} finds only nearby nodes.

    \item \textbf{\code{find\_lowest\_unvisited()}: $O(N)$ vs $O(\log N)$.} Original: iterates all nodes, finds minimum $z$ among unvisited. Refactored: uses \code{heapq} min-heap sorted by $z$.

    \item \textbf{DFS neighbor ordering.} Original: explores neighbors in \code{sorted(graph[node])} order with full recursion. Refactored: pushes neighbors onto stack in \code{reversed(sorted())} order. Can produce different pass compositions.

    \item \textbf{New: Strategy pattern.} \code{PathfindingStrategy} ABC with \code{generate\_print\_passes()} method. DFS and sweep-line are interchangeable via \code{--algorithm} CLI flag.
\end{enumerate}

\subsubsection{Performance on Test Network (31,271 points)}

\begin{table}[h]
\centering
\begin{tabular}{l l l}
\toprule
\textbf{Metric} & \textbf{Original} & \textbf{Refactored} \\
\midrule
Collision check    & $O(N)$ per node   & $O(k \log N)$ per node \\
Lowest unvisited   & $O(N)$ per pass   & $O(\log N)$ per pass \\
Recursion          & Crashes at default limit & No recursion \\
Total DFS time     & $\sim$4--5 minutes & $<$1 second \\
\bottomrule
\end{tabular}
\caption{Pathfinding performance comparison.}
\end{table}

%-----------------------------------------------------------------------------
\subsection{Pass Subdivision}
%-----------------------------------------------------------------------------

\textbf{Original:} Lines 997--1105 (109 lines)\\
\textbf{Refactored:} \code{core/postprocessing.py:11-84} (\code{subdivide\_passes}, 74 lines)\\
\textbf{Reduction:} $-$35 lines ($-$32.1\%)

Splits passes at DFS backtrack points; ensures bottom-up ordering.

\textbf{What changed:} Logic is identical. Reduction comes from removing changelog file writes and more concise Python idioms (e.g., \code{print\_passes[i].index(bp)} instead of manual iteration).

%-----------------------------------------------------------------------------
\subsection{Gap Closure Pipeline}
%-----------------------------------------------------------------------------

\textbf{Original:} Lines 1107--3200 (2,094 lines)\\
\textbf{Refactored:} \code{core/gap\_closure.py} (951 lines)\\
\textbf{Reduction:} $-$1,143 lines (\textbf{$-$54.6\%})

\textbf{This is the single largest reduction in the refactoring.}

6-part iterative gap closure: disconnect detection, condition 0 (vessel gaps), conditions 1--3 (branchpoint connections), final closure.

\subsubsection{Detailed Sub-section Mapping}

\begin{table}[h]
\centering
\footnotesize
\begin{tabularx}{\textwidth}{L{1.5cm} L{2.5cm} X R{0.6cm} R{0.6cm}}
\toprule
\textbf{Old Lines} & \textbf{Section} & \textbf{Refactored Function} & \textbf{Old} & \textbf{New} \\
\midrule
1107--1340 & Changelog 0 (disconnect detect) & \code{find\_disconnects()} (shared) & 234 & 244 \\
1341--1532 & Condition 0 (vessel gaps) & \code{close\_gaps\_condition0()} & 192 & 104 \\
1534--1758 & Changelog 1 (re-detect) & \code{find\_disconnects()} (reused) & 225 & 0 \\
1776--1906 & Branchpoint cond.\ 1 (backwards) & \code{close\_gaps\_branchpoint} \newline (dir=\code{"backwards"}) & 131 & 117 \\
1908--2137 & Changelog 2+3 (re-detect $\times$2) & \code{find\_disconnects()} (reused $\times$2) & 230 & 0 \\
2138--2259 & Branchpoint cond.\ 2 (forwards) & \code{close\_gaps\_branchpoint} \newline (dir=\code{"forwards"}) & 122 & 0 \\
2260--2488 & Changelog 4 (re-detect) & \code{find\_disconnects()} (reused) & 229 & 0 \\
2489--2572 & Branchpoint cond.\ 3 (parent) & \code{close\_gaps\_branchpoint} \newline (dir=\code{"parent\_to\_nbr"}) & 84 & 55 \\
2573--2802 & Changelog 5 (re-detect) & \code{find\_disconnects()} (reused) & 230 & 0 \\
2803--3099 & Final gap closure & \code{close\_gaps\_final()} + \newline \code{remove\_artifact\_passes()} & 297 & 107 \\
3100--3200 & Final cleanup & \code{run\_full\_gap\_closure\_pipeline()} & 101 & 151 \\
\bottomrule
\end{tabularx}
\caption{Gap closure pipeline detailed mapping.}
\end{table}

\subsubsection{How the 54.6\% Reduction Was Achieved}

\begin{enumerate}
    \item \textbf{Disconnect detection: 6 copies $\to$ 1 function.} The original copied the $\sim$230-line disconnect detection block \textbf{6 times} (changelogs 0--5), each nearly identical. The refactored code has a single \code{find\_disconnects()} function (244 lines) called 6 times. This alone saves $\sim$\textbf{1,150 lines}.

    \item \textbf{Branchpoint conditions: 3 similar blocks $\to$ 1 parameterized function.} Conditions 1, 2, and 3 had nearly identical structure but differed in search direction. The refactored \code{close\_gaps\_branchpoint()} takes a \code{direction} parameter (\code{"backwards"}, \code{"forwards"}, \code{"parent\_to\_neighbor"}).

    \item \textbf{Pipeline orchestrator.} \code{run\_full\_gap\_closure\_pipeline()} (151 lines) replaces the implicit sequencing of inline blocks with an explicit 12-step pipeline with logging at each stage.
\end{enumerate}

%-----------------------------------------------------------------------------
\subsection{Downsampling}
%-----------------------------------------------------------------------------

\textbf{Original:} Lines 3216--3310 (95 lines)\\
\textbf{Refactored:} \code{core/postprocessing.py:87-129} (\code{downsample\_passes}, 43 lines)\\
\textbf{Reduction:} $-$52 lines ($-$54.7\%)

Keeps every $N$th node while preserving first/last, branchpoints, endpoints.

\textbf{Reduction source:} The original included an inline plot generation block ($\sim$50 lines) that was moved to \code{\_write\_plots()} in the pipeline.

%-----------------------------------------------------------------------------
\subsection{Multimaterial Processing}
%-----------------------------------------------------------------------------

\textbf{Original:} Lines 3320--3695 (376 lines)\\
\textbf{Refactored:} \code{core/multimaterial.py} (98 lines) + \code{pipeline.py:352-434} (83 lines) --- total 181 lines\\
\textbf{Reduction:} $-$195 lines ($-$51.9\%)

Classifies passes by artven type, swaps noise, subdivides at material boundaries, computes speeds, generates MM plot.

\subsubsection{Detailed Mapping}

\begin{table}[h]
\centering
\footnotesize
\begin{tabularx}{\textwidth}{L{1.5cm} X L{3.8cm} R{0.6cm} R{0.6cm}}
\toprule
\textbf{Old Lines} & \textbf{What It Does} & \textbf{Refactored Location} & \textbf{Old} & \textbf{New} \\
\midrule
3320--3340 & Extract artven values per pass & Part of \code{\_subdivide\_by\_material()} & 21 & incl. \\
3357--3432 & Noise swapping (outlier artven at boundaries) & \code{\_subdivide\_by\_material()} \newline pipeline.py:354--371 & 76 & 18 \\
3433--3583 & Find break points, subdivide, reassemble & \code{\_subdivide\_by\_material()} \newline pipeline.py:372--434 & 151 & 65 \\
3610--3620 & Classify passes (last node's artven) & \code{classify\_passes\_\allowbreak by\_material()} \newline multimaterial.py:11--41 & 11 & 31 \\
3629--3695 & MM plot generation & \code{\_write\_plots()} + \newline \code{create\_network\_plot()} & 67 & 1 \\
3697--3727 & Speed: \code{flow / ($\pi r^2$)} & \code{compute\_radius\_speeds()} \newline multimaterial.py:44--75 & 31 & 32 \\
\bottomrule
\end{tabularx}
\caption{Multimaterial processing detailed mapping.}
\end{table}

\textbf{Reduction sources:} MM plot block (67 lines) replaced by 1-line call to parameterized function. Artven extraction simplified (no separate dictionary). Subdivision logic more concise with Python idioms.

%-----------------------------------------------------------------------------
\subsection{Overlap (SM and MM)}
%-----------------------------------------------------------------------------

\textbf{Original:} SM: Lines 3729--3806 (78 lines); MM: Lines 3880--3959 (80 lines)\\
\textbf{Refactored:} \code{core/postprocessing.py:132-242} (111 lines total)\\
\textbf{Reduction:} $-$47 lines

Finds passes ending on shared nodes in earlier passes, retraces backwards to add overlap nodes.

\textbf{What changed:}

\begin{enumerate}
    \item \textbf{SM and MM overlap: 2 near-identical blocks $\to$ 1 function.} The original had separate SM and MM blocks with the same logic. The refactored \code{add\_overlap()} works on any pass dictionary.

    \item \textbf{Two algorithms available.} The \code{retrace} algorithm matches the original. The \code{consecutive} algorithm is a new faster alternative that only checks adjacent pass pairs.
\end{enumerate}

%-----------------------------------------------------------------------------
\subsection{Custom Gap Closure Files}
%-----------------------------------------------------------------------------

\textbf{Original:} Lines 3808--3876 (69 lines)\\
\textbf{Refactored:} \code{pipeline.py:437-462} (\code{\_load\_gap\_extensions}, 26 lines)\\
\textbf{Reduction:} $-$43 lines ($-$62.3\%)

Reads pass indices and delta vectors from extension files.

\textbf{Reduction source:} The original had separate SM and MM reading blocks. The refactored code has a single function called with a suffix parameter (\code{"SM"} or \code{"MM"}).

%-----------------------------------------------------------------------------
\subsection{Coordinate Output Files}
%-----------------------------------------------------------------------------

\textbf{Original:} SM: Lines 4090--4251 (162 lines); MM: Lines 4253--4405 (153 lines)\\
\textbf{Refactored:} \code{io/writer.py} (148 lines)\\
\textbf{Reduction:} $-$167 lines ($-$53.0\%)

Writes x, y, z, radius, speed, and combined coordinate files per pass.

\textbf{How the reduction was achieved:}

\begin{enumerate}
    \item \textbf{SM and MM: 2 blocks $\to$ 1 parameterized function.} \code{write\_pass\_coordinates()} handles both SM and MM by parameter.

    \item \textbf{Per-column output: repeated inline loops $\to$ \code{\_write\_column\_file()} helper.} The original had separate loop blocks for x, y, z, radius, and artven output.

    \item \textbf{Speed output separated.} \code{\_write\_speed\_file()} handles the different speed dictionary format.
\end{enumerate}

%-----------------------------------------------------------------------------
\subsection{G-code Generation}
%-----------------------------------------------------------------------------

\textbf{Original:} Lines 4409--5486 (1,078 lines across 6 blocks)\\
\textbf{Refactored:} \code{io/gcode/} package (713 lines across 4 files)\\
\textbf{Reduction:} $-$365 lines (\textbf{$-$33.9\%})

\textbf{This is the second largest reduction in the refactoring.}

\subsubsection{Detailed Block Mapping}

\begin{table}[h]
\centering
\footnotesize
\begin{tabularx}{\textwidth}{X R{0.7cm} L{3.8cm} R{0.7cm} R{0.8cm}}
\toprule
\textbf{Original Block} & \textbf{Old} & \textbf{Refactored Class} & \textbf{New} & \textbf{$\Delta$} \\
\midrule
SM Pressure (4409--4506)      & 98  & \code{PressureGcodeWriter}    & $\sim$70  & $-$28 \\
SM Positive Ink (4508--4629)  & 122 & \code{PositiveInkGcodeWriter} & $\sim$85  & $-$37 \\
SM Aerotech (4632--4711)      & 80  & \code{AerotechGcodeWriter}    & $\sim$60  & $-$20 \\
MM Pressure (4713--4974)      & 262 & \code{PressureGcodeWriter}    & $\sim$70  & $-$192 \\
MM Positive Ink (4976--5277)  & 302 & \code{PositiveInkGcodeWriter} & $\sim$95  & $-$207 \\
MM Aerotech (5279--5486)      & 208 & \code{AerotechGcodeWriter}    & $\sim$90  & $-$118 \\
(not in original)             & 0   & \code{GcodeWriter} ABC        & 183       & +183 \\
\midrule
\textbf{Total}                & \textbf{1,072} & & \textbf{713} & \textbf{$-$365} \\
\bottomrule
\end{tabularx}
\caption{G-code generation block mapping. SM/MM pairs share one class each.}
\end{table}

\subsubsection{How the Reduction Was Achieved}

\begin{enumerate}
    \item \textbf{Abstract base class \code{GcodeWriter} (183 lines).} Extracts the shared iteration loop (header $\to$ for each pass $\to$ for each node $\to$ footer), gap extension logic, and file I/O. All 6 original blocks had this same loop structure copy-pasted.

    \item \textbf{SM/MM merged into single classes.} Each original pair (SM + MM) was merged into one class with \code{if config.multimaterial:} branches.

    \item \textbf{Custom G-code loading: per-pass file open $\to$ load-once.} Original: opened and read the same file on \textbf{every pass start} (potentially 40+ times). Refactored: \code{CustomCodes.load\_from\_dir()} reads all 12 template files once into memory at startup.

    \item \textbf{Shared methods eliminate duplication:}
    \begin{itemize}
        \item \code{\_write\_header()} --- 1 implementation instead of 6 inline blocks
        \item \code{\_write\_footer()} --- 1 implementation instead of 6 inline blocks
        \item \code{\_write\_network()} --- 1 shared loop instead of 6 copy-pasted loops
        \item \code{\_write\_custom()} --- 1 helper instead of repeated \code{with open(...) as f:}
    \end{itemize}
\end{enumerate}

%-----------------------------------------------------------------------------
\subsection{Print Instructions}
%-----------------------------------------------------------------------------

\textbf{Original:} Lines 5488--5568 (81 lines)\\
\textbf{Refactored:} \code{pipeline.py:516-619} (\code{\_print\_instructions}, 104 lines)

Computes bounding box, centering positions, prints SM/MM operator instructions.

\textbf{What changed:} Logic is identical. Line increase due to function signature, docstring, and string formatting across multiple lines for readability. The instructions themselves are character-for-character identical (validated by test output comparison).

\textbf{Note:} This function was \textbf{completely missing} from the initial refactoring and was added during the parity fixes.

%-----------------------------------------------------------------------------
\subsection{Plot Generation (All Variants)}
%-----------------------------------------------------------------------------

\textbf{Original:} 5--7 inline blocks totaling $\sim$335 lines\\
\textbf{Refactored:} \code{viz/plotting.py} (137 lines) + calls in \code{pipeline.\_write\_plots()} (49 lines) --- total 186 lines\\
\textbf{Reduction:} $-$149 lines ($-$44.5\%)

\subsubsection{Original Plot Blocks}

\begin{table}[h]
\centering
\small
\begin{tabularx}{\textwidth}{l l X}
\toprule
\textbf{Lines} & \textbf{Plot} & \textbf{Refactored} \\
\midrule
328--394 (67)    & Original network          & \code{create\_original\_network\_plot()} (60 lines) \\
3100--3155 (56)  & Final SM passes           & \code{create\_network\_plot()} (62 lines, shared) \\
3258--3310 (53)  & Downsampled SM            & \code{create\_network\_plot()} (reused) \\
3629--3695 (67)  & MM passes (art./venous)   & \code{create\_network\_plot(colors=...)} (reused) \\
3965--4018 (54)  & SM overlap                & \code{create\_network\_plot()} (reused) \\
\bottomrule
\end{tabularx}
\caption{Original plot blocks and their refactored equivalents.}
\end{table}

\textbf{Reduction mechanism:} 5 copy-pasted plot blocks replaced by 2 parameterized functions. Each original block was $\sim$55--67 lines of nearly identical Plotly code with minor variations (title, color, data source).

%=============================================================================
\section{New Features (Not in Original)}
%=============================================================================

\begin{table}[h]
\centering
\small
\begin{tabular}{l l r}
\toprule
\textbf{Feature} & \textbf{Location} & \textbf{Lines} \\
\midrule
Streamlit GUI                     & \code{gui/app.py}              & 792 \\
Sweep-line pathfinding            & \code{core/pathfinding.py}     & 147 \\
Nearest-neighbor pass reordering  & \code{core/postprocessing.py}  & 46 \\
Spatial index module              & \code{spatial/index.py}        & 120 \\
Test suite                        & \code{tests/} (5 files)        & 501 \\
\code{XcavateConfig} dataclass    & \code{config.py}               & 131 \\
Progress callbacks                & \code{pipeline.py}             & $\sim$10 \\
Structured logging                & Throughout                     & $\sim$20 \\
\bottomrule
\end{tabular}
\caption{New features added during refactoring.}
\end{table}

\textbf{Total new code not derived from original: $\sim$1,767 lines.} This accounts for the overall increase from 5,568 to 6,201 lines despite significant reductions in existing code.

%=============================================================================
\section{Reduction Summary}
%=============================================================================

\begin{table}[h]
\centering
\small
\begin{tabularx}{\textwidth}{X r r r r}
\toprule
\textbf{Section} & \textbf{Original} & \textbf{Refact.} & \textbf{Change} & \textbf{\%} \\
\midrule
Gap closure pipeline        & 2,094  & 951   & $-$1,143 & $-$54.6\% \\
G-code generation (6 blks)  & 1,078  & 713   & $-$365   & $-$33.9\% \\
Plot generation (5--7 blks)  & $\sim$335 & 186 & $-$149   & $-$44.5\% \\
Coordinate output (SM+MM)   & 315    & 148   & $-$167   & $-$53.0\% \\
Overlap (SM+MM)             & 158    & 111   & $-$47    & $-$29.7\% \\
Downsampling + plot         & 95     & 43    & $-$52    & $-$54.7\% \\
Multimaterial processing    & 376    & 181   & $-$195   & $-$51.9\% \\
Gap extension file loading  & 69     & 26    & $-$43    & $-$62.3\% \\
File reading/preprocessing  & 148    & 298   & +150     & +101.4\% \\
Graph construction          & 418    & 836   & +418     & +100.0\% \\
DFS pathfinding             & 85     & 620   & +535     & +629.4\% \\
Subdivision                 & 109    & 74    & $-$35    & $-$32.1\% \\
Print instructions          & 81     & 104   & +23      & +28.4\% \\
Args/config/setup           & 174    & 332   & +158     & +90.8\% \\
\midrule
\textbf{Sub-total (ported code)} & \textbf{5,535} & \textbf{4,623} & \textbf{$-$912} & \textbf{$-$16.5\%} \\
New features (GUI, tests, \ldots) & 0 & 1,578 & +1,578 & N/A \\
\midrule
\textbf{Grand total}        & \textbf{5,568} & \textbf{6,201} & \textbf{+633} & \textbf{+11.4\%} \\
\bottomrule
\end{tabularx}
\caption{Reduction summary across all sections.}
\end{table}

\textbf{Key insight:} The ported code is 16.5\% shorter despite being more modular, better documented, and having more robust algorithms. The overall 11.4\% increase is entirely due to new features (GUI, sweep-line, tests, spatial indexing, config system).

%=============================================================================
\section{Algorithmic Complexity Improvements}
%=============================================================================

\begin{table}[h]
\centering
\small
\begin{tabularx}{\textwidth}{L{3.5cm} X X}
\toprule
\textbf{Operation} & \textbf{Original} & \textbf{Refactored} \\
\midrule
Collision detection (\code{is\_valid})    & $O(N)$ per call            & $O(k \log N)$ per call \\
Lowest unvisited node                    & $O(N)$ per pass            & $O(\log N)$ amortized \\
Interpolation ($N$ insertions)           & $O(N^2)$ total copies      & $O(N)$ single concatenation \\
Branchpoint detection                    & $O(E \times N)$ nested loops & $O(E \times k \times \log N)$ KD-tree \\
Inlet/outlet matching                    & $O(N \times I)$ exact cmp   & $O(I \times \log N)$ KD-tree \\
Coordinate rounding                      & $O(N \times C)$ Python loop & $O(1)$ vectorized numpy \\
Custom G-code reading                    & $O(P \times F)$ file opens  & $O(F)$ load once \\
\bottomrule
\end{tabularx}
\caption{Algorithmic complexity improvements. $N \approx 31{,}000$ (points), $E \approx 22$ (endpoints), $I = 2$ (inlets/outlets), $P \approx 45$ (passes), $F = 12$ (template files), $C = 3$--$5$ (columns), $k \approx 5$--$10$ (nearby neighbors).}
\end{table}

%=============================================================================
\section{Design Pattern Changes}
%=============================================================================

\begin{table}[h]
\centering
\small
\begin{tabularx}{\textwidth}{L{3cm} X X}
\toprule
\textbf{Pattern} & \textbf{Original} & \textbf{Refactored} \\
\midrule
Configuration       & 30+ global variables                & \code{XcavateConfig} dataclass with typed fields \\
G-code generation   & 6 copy-pasted blocks with inline conditionals & ABC + 3 subclasses (Template Method pattern) \\
Pathfinding         & Single recursive function            & Strategy pattern with pluggable algorithms \\
Plot generation     & 5--7 inline Plotly blocks            & 2 parameterized functions (DRY) \\
File I/O            & Inline \code{open/read/write} throughout & Dedicated reader/writer modules \\
Coordinate output   & 2 copy-pasted blocks (SM + MM)       & 1 parameterized function \\
Disconnect detection & 6 copy-pasted blocks                & 1 function called 6 times \\
Branchpoint conds   & 3 similar blocks                    & 1 parameterized function with direction parameter \\
Spatial queries     & Brute-force distance computation     & KD-tree spatial indexing \\
Error handling      & None (crashes on bad input)           & Validation, warnings, graceful errors \\
Progress reporting  & \code{print()} statements            & \code{logging} module + GUI progress callbacks \\
Entry points        & Script execution only                & CLI + GUI + library import \\
\bottomrule
\end{tabularx}
\caption{Design pattern changes from original to refactored code.}
\end{table}

\end{document}
